% === Begin Label Data ===
% ID: 676
% 难度星级: 
% 标签: 
% 备注: 
% 组卷引用次数: 0
% === End  Label Data ===

\begin{problem}{2008}{G}{湖南卷}{15}{概率}
对有 $n$ ($n\geqslant 4$) 个元素的总体 $\{1,2,3,\cdots,n\}$ 进行抽样，先将总体分成两个子总体 $\{1,2,\cdots,m\}$ 和 $\{m+1,m+2,\cdots,n\}$ ($m$ 是给定的正整数，且 $2\leqslant m\leqslant n-2$)，再从每个子总体中各随机抽取 $2$ 个元素组成样本，用 $P_{ij}$ 表示元素 $i$ 和 $j$ 同时出现在样本中的概率，则 $P_{1n}=$ \underline{\hspace{4em}}；所有 $P_{ij}$ ($1\leqslant i<j\leqslant n$) 的和等于 \underline{\hspace{4em}}.
\end{problem}

\begin{answer}
$ \dfrac{4}{m(n-m)} $；$ 6 $
\end{answer}

\begin{solutions}
（1）由题意有，$ P_{1n}=\dfrac{\mathrm{C}_{m-1}^{1}\mathrm{C}_{n-m-1}^{1}}{\mathrm{C}_{m}^{2}\mathrm{C}_{n-m}^{2}}=\dfrac{4}{m(n-m)} $．

（2）当 $ 1\leqslant i<j\leqslant m $ 时，$ P_{ij}=\dfrac{1}{\mathrm{C}_{m}^{2}} $，这样的 $ P_{ij} $ 有 $ \mathrm{C}_{m}^{2} $ 个，  
故所有 $ P_{ij}\,(1\leqslant i<j\leqslant m) $ 的和为 $ \dfrac{1}{\mathrm{C}_{m}^{2}}\cdot\mathrm{C}_{m}^{2}=1 $；  

当 $ m+1\leqslant i<j\leqslant n $ 时，$ P_{ij}=\dfrac{1}{\mathrm{C}_{n-m}^{2}} $，这样的 $ P_{ij} $ 有 $ \mathrm{C}_{n-m}^{2} $ 个，  
故所有 $ P_{ij}\,(m+1\leqslant i<j\leqslant n) $ 的和为 $ \dfrac{1}{\mathrm{C}_{n-m}^{2}}\cdot\mathrm{C}_{n-m}^{2}=1 $；  

当 $ 1\leqslant i\leqslant m $，$ m+1\leqslant j\leqslant n $ 时，$ P_{ij}=\dfrac{1}{m(n-m)} $，这样的 $ P_{ij} $ 共有 $ m(n-m) $ 个，  
故所有 $ P_{ij}\,(1\leqslant i\leqslant m,\,m+1\leqslant j\leqslant n) $ 的和为 $ \dfrac{1}{m(n-m)}\cdot m(n-m)=4 $．  

综上所述，所有 $ P_{ij}\,(1\leqslant i<j\leqslant n) $ 的和等于 $ 1+1+4=6 $．
\end{solutions}